\documentclass[../../../include/open-logic-section]{subfiles}

\begin{document}

\olfileid{sth}{replacement}{absinf}
\olsection{الاستبدال و«اللانهاية المطلقة»}

ندع الآن~\limofsize{} جانبًا، وننظر في طائفة أخرى من المحاولات
الداخلية لتسويغ الاستبدال؛ وهي محاولات تحمل فكرة تكون المجموعات على
مراحل محمل الجد.

حين رسمنا العملية التكرارية أول مرة، وضعنا مبادئ تبين ما يجري في كل
مرحلة، وهي~\stageshier و\stagesord و\stagesacc. ثم زدنا مبادئ تتعلق
بعدد المراحل: فـ\stagessucc{} يقضي بألا تنتهي عملية تكوين المجموعات،
و\stagesinf{} يقضي بأن تمر العملية بمرحلة لا نهائية.

وقد يقال إن المبدأين الأخيرين فرعان عن أصل أعم، نحو:
\begin{enumerate}
	\item[] \stagesinex. عدد المراحل لا نهائي على نحو مطلق؛ والهرم
	مرتفع إلى أقصى ارتفاع ممكن.
\end{enumerate}
وهذا، ولا ريب، مبدأ غير صوري. لكنه، وإن لم يلزم \emph{مباشرة}
من التصور التراكمي-التكراري، يبدو \emph{موافقًا} له. فهو، بخلاف
\limofsize{}، يبقي على معنى تكون المجموعات مرحلة إثر مرحلة.

فنرجو أن نتخذ~\stagesinex{} سبيلًا إلى تسويغ الاستبدال، ولننظر كيف
يكون ذلك.

مر بنا في~\olref[ordinals][idea]{sec} أن من اليسير إنشاء ترتيب جيد
ينبغي، بحسب المقصود، أن يتماثل مع~$\omega+\omega$. أي إننا نقدر على
تصور عملية تكرارية تسير مرحلة بعد مرحلة ويكون نمط ترتيبها، بحسب
المقصود،~$\omega+\omega$. فإذا سلمنا~\stagesinex، وجب أن نسلم بوجود
مرحلة الهرم ذات الرتبة~$\omega+\omega$ على الأقل، وهي
$\OLId{latin-looped}{V}_{\omega+\omega}$؛ إذ لا شك في أن الهرم \emph{يقدر} أن يمتد إلى
هذا الحد.

وهذا الوجه يعم كل ترتيب جيد: ينبغي لعملية بناء الهرم التكراري أن
تمتد بقدر ذلك الترتيب على الأقل. ويكفينا لضمانه أن نتخذ
\olref[ordinals][ordtype]{thmOrdinalRepresentation} \emph{مسلّمة}؛
فهي تقضي بأن كل ترتيب جيد يتماثل مع عدد ترتيبي لفون نويمان. ولما كان
كل عدد ترتيبي لفون نويمان مساويًا لرتبته، لزم من
\olref[ordinals][ordtype]{thmOrdinalRepresentation} أن كل ترتيب جيد
يمكن وصفه في نظريتنا تسبقه عملية البناء التكراري للمجموعات.

وهذا شبيه بأن يكون نتيجة لـ\stagesinex. غير أن بينه وبين استخراج
الاستبدال عقبة فنية: فالاستبدال أقوى تمامًا من
\olref[ordinals][ordtype]{thmOrdinalRepresentation}.
(وقد نبه~\citet[\S13.2]{Potter2004} إلى ذلك، وسيأتي برهانه في
\olref[sth][replacement][finiteaxiomatizability]{sec}.)

فإذا أردنا أن نفهم~\stagesinex{} فهمًا يستتبع الاستبدال، لم يجز أن
يقتصر معناه على أن الهرم يجاوز كل ترتيب جيد؛ بل لا بد من معنى أقوى.
ولهذا اقترح~\citet{Shoenfield:AST} تقوية طبيعية، وهي أن الهرم لا
\emph{تسايره} مجموعة.
\footnote{يبدو أن غودل اقترح فكرة مشابهة؛ انظر
  \citet[p.~223]{Potter2004}. ولمناقشة غودل و
  \citeauthor{Shoenfield:AST}، انظر~\citet[90--5]{Incurvati2020}.}
وتفصيل ذلك: إذا كان~$\tau$ تطبيقًا يرسل المجموعات إلى مراحل من الهرم،
لم تستنفد صورة أي مجموعة~$\OLId{latin-looped}{A}$ تحت~$\tau$ الهرم. وصيغته المخططية:
\begin{enumerate}
	\item[] \stagescofin. إذا كانت~$\OLId{latin-looped}{A}$ مجموعة وكانت~$\tau(\OLId{latin-ordinary}{x})$ مرحلة
	لكل~$\OLId{latin-ordinary}{x} \in \OLId{latin-looped}{A}$، فهناك مرحلة تأتي بعد كل~$\tau(\OLId{latin-ordinary}{x})$، حيث~$\OLId{latin-ordinary}{x} \in \OLId{latin-looped}{A}$.
\end{enumerate}
ومن البين أن~$\ZF$ تثبت صيغة ملائمة من~\stagescofin. وفي الجهة
الأخرى نستطيع أن نبين، ببرهان غير صوري، أن~\stagescofin{} يسوّغ
الاستبدال.
\footnote{سيكون إثبات الاستبدال باستعمال صياغة صورية ما لـ
  \stagescofin أصعب، لأن~$\Z$ وحدها ليست قوية بما يكفي لتعريف
  المراحل، ولذلك لا يتضح كيف يمكن صوغ~\stagescofin صوريًا. لكن أحد
  الخيارات هو العمل في امتداد ما لـ$\LT$، كما نوقش في
  \olref[sth][replacement][extrinsic]{sec}.}
افترض إذن~$(\forall \OLId{latin-ordinary}{x} \in \OLId{latin-looped}{A})\lexists![\OLId{latin-ordinary}{y}][\phi(\OLId{latin-ordinary}{x},\OLId{latin-ordinary}{y})]$. ولكل~$\OLId{latin-ordinary}{x} \in \OLId{latin-looped}{A}$،
لتكن~$\sigma(\OLId{latin-ordinary}{x})$ هي~$\OLId{latin-ordinary}{y}$ التي تحقق~$\phi(\OLId{latin-ordinary}{x},\OLId{latin-ordinary}{y})$، ولتكن~$\tau(\OLId{latin-ordinary}{x})$ أول
مرحلة تتكون فيها~$\sigma(\OLId{latin-ordinary}{x})$. يضمن~\stagescofin{} مرحلة~$\OLId{latin-looped}{V}$ بحيث
$(\forall \OLId{latin-ordinary}{x} \in \OLId{latin-looped}{A})\tau(\OLId{latin-ordinary}{x})\in \OLId{latin-looped}{V}$. ولما كان كل~$\tau(\OLId{latin-ordinary}{x}) \in \OLId{latin-looped}{V}$،
وكانت~$\sigma(\OLId{latin-ordinary}{x}) \subseteq \tau(\OLId{latin-ordinary}{x})$، وكانت~$\OLId{latin-looped}{V}$ مقتدرة، حصل
$\sigma(\OLId{latin-ordinary}{x}) \in \OLId{latin-looped}{V}$. ثم يعطينا الفصل
$\Setabs{\OLId{latin-ordinary}{y} \in \OLId{latin-looped}{V}}{(\exists \OLId{latin-ordinary}{x} \in \OLId{latin-looped}{A})\sigma(\OLId{latin-ordinary}{x}) = \OLId{latin-ordinary}{y}} =
\Setabs{\OLId{latin-ordinary}{y}}{(\exists \OLId{latin-ordinary}{x} \in \OLId{latin-looped}{A})\phi(\OLId{latin-ordinary}{x},\OLId{latin-ordinary}{y})}$.

\begin{prob}
	صُغ~\stagescofin{} صوريًا داخل~$\ZF$.
\end{prob}

فقد تبين أن~\stagescofin{} يسوّغ الاستبدال، ومن المعقول على الأقل أن
يكون~\stagescofin{} لازمًا عن~\stagesinex. فلنفرض إخفاق
\stagescofin؛ حينئذ تساير الهرم مجموعة~$\OLId{latin-looped}{A}$، أي يوجد تطبيق~$\tau$
بحيث لكل مرحلة~$\OLId{latin-looped}{S}$ بعض~$\OLId{latin-ordinary}{x} \in \OLId{latin-looped}{A}$ يحقق~$\OLId{latin-looped}{S} \in \tau(\OLId{latin-ordinary}{x})$. فتكون لنا
وسيلة للإحاطة بـ«اللانهاية المطلقة» المزعومة للهرم، إذ \emph{يستنفدها}
مدى~$\tau$ على~$\OLId{latin-looped}{A}$. وهذا ينافي قوة القول إن الهرم «لا نهائي على نحو
مطلق». فبالمقابل بالنقيض يستلزم~\stagesinex{} المبدأ~\stagescofin،
وهو بدوره يسوّغ الاستبدال.

فهذه محاولة جديرة بالاعتبار في تسويغ الاستبدال من داخله، ويبقى إليك
الحكم في تمام نجاحها.

\end{document}
